\subsection{Orientation of a graph}

Let G be a graph. An orientation of G is a subset A of the set of arrows of G such that \(A \cap \overline{A} = \varnothing\) and \(A \cup \overline{A} = \mathrm{Fl}(G)\) .

A graph equipped with an orientation is called an oriented graph.

Let G be an oriented graph and let A be its orientation; an oriented subgraph of G is a subgraph \(G'\) of G equipped with the orientation \(A' = \mathrm{Fl}(G') \cap A\) .

Let (G, A) and (G', A') be oriented graphs. A morphism of oriented graphs from (G, A) to (G', A') is a graph morphism \(\varphi\colon G \to G'\) such that \(\operatorname{Fl}(\varphi)(A) \subset A'\) .

